\documentclass[11pt]{article}
\usepackage[margin=0.85in]{geometry}
\usepackage[T1]{fontenc}
\usepackage{lmodern,amsmath,listings,booktabs,tabularx,graphicx,float}
\usepackage[hidelinks]{hyperref}
\hypersetup{pdftitle={Assignment 2, Part 3(a): CUBIC vs. Reno vs. BBR},
  pdfauthor={CS422 Fall 2026 project team}}
\lstset{basicstyle=\ttfamily\small,columns=fullflexible,keepspaces=true,
  breaklines=true,frame=single,showstringspaces=false}
\setlength{\parindent}{0pt}
\setlength{\parskip}{0.45em}
% Long code names cannot hyphenate; let lines stretch instead of overflowing.
\sloppy
\setlength{\emergencystretch}{3em}

% Permalinks to the exact code that produced these results.
\newcommand{\repo}{https://github.com/LelsersLasers/CS422-Fall2026/blob/ea2af041267d28c0d6a68aca790cdb8e741e6938/Assignment2}
\newcommand{\results}{part3a-results/comparison.pdf}

\begin{document}
\begin{center}
{\Large Assignment 2, Part 3(a): CUBIC vs.\ Reno vs.\ BBR}\\[0.5em]
{\small Ayusha Patra \quad Kabileswaran Krishna Kumar \quad
Vihaan Kodeboyina \quad Millan Kumar}
\end{center}

\section*{Summary}

Across five public iperf3 servers and three repetitions, \textbf{BBR had the
highest goodput in all 15 comparisons} (78.1\,Mbit/s mean vs.\ 52.1 for CUBIC
and 43.6 for Reno), but added 12--19\,ms of queueing delay and retransmitted
about 11\% of its bytes. CUBIC beat Reno in 13 of 14 comparisons, by more on
longer paths. The paths were \emph{loss-limited}: sparse losses kept CUBIC and
Reno at a roughly constant window, while BBR, which ignores loss, was not
limited this way.

\section*{Setup}

\begin{tabularx}{\textwidth}{@{}lX@{}}
\toprule
Host & Ubuntu 24.04.5 VM (Multipass) on an Apple Silicon Mac, Wi-Fi uplink \\
Kernel & Linux \texttt{6.8.0-147-generic} (aarch64), BBRv1 via
  \texttt{modprobe tcp\_bbr} \\
Command & \texttt{python3 main.py -n 5 -d 60 -i 0.2 -r 3 -{}-seed 42} \\
Tests & 60\,s each, TCP\_INFO sampled every 0.2\,s, 3 repetitions
  $\times$ 5 destinations $\times$ 3 algorithms = 45 \\
\bottomrule
\end{tabularx}

Part 3 reuses the five destinations from Part 1 (minimum RTT in
parentheses): LeaseWeb US (48\,ms), FortiSASE ES (105\,ms), MilkyWan FR
(109\,ms), Multinet24 PL (130\,ms), and Mauritius Telecom (288\,ms). In each
repetition, every destination runs all three algorithms back to back, and the
order rotates: CUBIC$\rightarrow$Reno$\rightarrow$BBR, then
Reno$\rightarrow$BBR$\rightarrow$CUBIC, then BBR$\rightarrow$CUBIC$\rightarrow$Reno.
Each test is preceded by a 3\,s pause, and refused tests are retried up to
twice. One test (Reno, repetition 3, LeaseWeb) was refused every time, so
44 of 45 tests completed. The algorithm is set with
\texttt{setsockopt(TCP\_CONGESTION)} before connecting and verified with
\texttt{getsockopt}.

\textbf{Metrics} (all from TCP\_INFO on the data socket): goodput from
$\Delta$\texttt{bytes\_acked}; smoothed RTT; queueing delay $=$ mean RTT
$-$ \texttt{min\_rtt}; loss proxy $=$ retransmitted bytes / bytes sent, and
retransmissions per interval in the plots.

\section*{Results}

\begin{table}[!htbp]
\centering
\small
\caption{Per destination: mean $\pm$ std.\ dev.\ of per-run goodput across
repetitions, and means of the other metrics.}
\label{tab:dest}
\begin{tabular}{@{}llrrrrr@{}}
\toprule
Destination & Algo & Goodput & RTT & Queueing & Retrans. & cwnd \\
(min RTT) & & (Mbit/s) & (ms) & (ms) & bytes & (segments) \\
\midrule
LeaseWeb US & CUBIC & 93.2 $\pm$ 1.4 & 77 & 28.7 & 0.04\% & 601 \\
(48 ms) & Reno & 85.7 $\pm$ 3.7 & 72 & 25.6 & 0.36\% & 538 \\
 & BBR & \textbf{144.5} $\pm$ 11.9 & 91 & 44.5 & 12.64\% & 2838 \\
\midrule
FortiSASE ES & CUBIC & 48.6 $\pm$ 1.4 & 133 & 25.6 & 0.07\% & 588 \\
(105 ms) & Reno & 45.0 $\pm$ 1.6 & 129 & 23.3 & 1.01\% & 544 \\
 & BBR & \textbf{67.5} $\pm$ 1.4 & 141 & 35.7 & 9.08\% & 1961 \\
\midrule
MilkyWan FR & CUBIC & 47.1 $\pm$ 2.0 & 132 & 24.1 & 0.05\% & 571 \\
(109 ms) & Reno & 45.0 $\pm$ 1.4 & 132 & 24.4 & 0.66\% & 546 \\
 & BBR & \textbf{72.8} $\pm$ 6.2 & 142 & 36.0 & 15.51\% & 2327 \\
\midrule
Multinet24 PL & CUBIC & 41.5 $\pm$ 3.9 & 153 & 22.3 & 0.52\% & 605 \\
(130 ms) & Reno & 36.0 $\pm$ 2.1 & 155 & 23.8 & 0.54\% & 529 \\
 & BBR & \textbf{63.3} $\pm$ 5.3 & 168 & 37.5 & 11.43\% & 2360 \\
\midrule
Mauritius Tel. & CUBIC & 30.0 $\pm$ 5.5 & 321 & 27.8 & 0.51\% & 834 \\
(288 ms) & Reno & 20.2 $\pm$ 2.5 & 316 & 22.4 & 1.00\% & 565 \\
 & BBR & \textbf{42.4} $\pm$ 1.9 & 332 & 39.3 & 7.39\% & 3207 \\
\bottomrule
\end{tabular}
\end{table}

The plots use \texttt{148.230.45.213:30009}, the representative destination
from Question 2, with shared axis scales for all three algorithms.

\begin{figure}[!htbp]
\centering
\includegraphics[page=1,width=\textwidth]{\results}
\caption{Time series at the representative destination (repetition 1).}
\label{fig:ts}
\end{figure}

\begin{figure}[!htbp]
\centering
\includegraphics[page=2,width=\textwidth]{\results}
\caption{cwnd, RTT, and retransmissions vs.\ goodput (all 0.2\,s samples,
all three repetitions).}
\label{fig:scatter}
\end{figure}

\begin{figure}[!htbp]
\centering
\includegraphics[page=9,width=0.9\textwidth]{\results}
\caption{Run-to-run variability, one point per test (all destinations).}
\label{fig:var}
\end{figure}

\section*{Analysis}

\textbf{Traces (Figure~\ref{fig:ts}).} Reno grows its window by about one
segment per RTT and halves it on loss (one loss in about 40\,s here). CUBIC
shows its cubic sawtooth, dropping to 70\% and regrowing within seconds. BBR
keeps a 1,300--3,300 segment window and dips almost to zero at about 12, 31,
and 48\,s; these are its ProbeRTT phases. BBR is the only algorithm with
sustained retransmissions (Figure~\ref{fig:scatter}).

\textbf{Operating regime: loss-limited.} On the four paths from 48 to
130\,ms, CUBIC and Reno kept nearly the same mean window (530--610 segments),
so goodput $\approx$ cwnd/RTT fell roughly as $1/\text{RTT}$. A capacity
limit would instead make the window track each path's BDP. A constant window
matches the AIMD throughput model, cwnd $\approx 1.22/\sqrt{p}$: Reno's one
loss per $\sim$155{,}000 segments ($p\approx 6\times10^{-6}$) predicts about
480 segments, close to the 530--565 measured. BBR does not back off on loss,
so these sparse losses did not limit it.

\textbf{Which algorithm is better where.}
\begin{itemize}
\item \textbf{Throughput: BBR}, by 1.5--1.8$\times$ Reno on the
  48--130\,ms paths and 2.1$\times$ at 288\,ms. On a long path, a loss-based
  algorithm needs more round trips to regrow after each backoff. BBR suits
  bulk transfer on high-BDP paths with sparse non-congestive loss.
\item \textbf{CUBIC vs.\ Reno: CUBIC on high-BDP paths.} The gap was
  1.05--1.15$\times$ up to 130\,ms but 1.48$\times$ at 288\,ms (window 834
  vs.\ 565). CUBIC backs off less (to 70\%, not 50\%), and its growth depends
  on time since the last loss rather than on RTT, whereas Reno would need
  over 2 minutes to regrow 500 segments at 288\,ms.
\item \textbf{Delay and loss: CUBIC and Reno.} Both had about 22--28\,ms of
  queueing delay on every path. Because this floor is the same everywhere,
  it is probably jitter from our Wi-Fi and VM networking. BBR added 12--19\,ms
  above it and retransmitted 7--16\% of its bytes: BBRv1 keeps up to
  2$\times$ the estimated BDP in flight and probes at 1.25$\times$ its
  bandwidth estimate, so it repeatedly overfills the bottleneck buffer.
\item \textbf{BDP.} The BDP implied by achieved goodput is about
  0.5--0.9\,MB on the shorter paths and up to 1.5\,MB at 288\,ms. CUBIC and
  Reno's $\sim$0.8\,MB windows are near it on short paths but well below it on
  the long one, where they trail BBR most. BBR's 2.8--4.6\,MB windows exceed
  the BDP, and the excess becomes queueing delay and loss.
\end{itemize}

We did not observe a deep-buffer regime, where BBR should have the
\emph{lower} RTT, or a short uncongested path, where all three should reach
link capacity.

\textbf{Variability.} On the same destination, results were consistent:
the coefficient of variation of per-run goodput was 1.5--9.4\% for 13 of 15
destination--algorithm pairs, and up to 18\% on the long Mauritius path
(Figure~\ref{fig:var}). Across destinations, goodput varied 3--4$\times$, so
the path matters more than the repetition. The ranking held regardless of
test order: BBR was fastest in all 15 groups, and CUBIC beat Reno in 13 of 14
(the exception was a 0.9\,Mbit/s difference).

\textbf{Limitations.} Public paths change between tests, and three
repetitions is a small sample. We tested one flow at a time, so we could not
measure fairness or BBR's effect on competing flows. The bottleneck location
and buffer size are unknown, and our VM on Wi-Fi adds jitter. BDP is a lower
bound because true capacity is unknown. TCP\_INFO sampling every 0.2\,s
averages short bursts. Results apply to BBRv1 only; later versions react to
loss.

\section*{Reproduction and code}

From \texttt{Assignment2/} on Linux with \texttt{reno cubic bbr} available:

\begin{lstlisting}
python3 main.py -n 5 -d 60 -i 0.2 -r 3 --seed 42
\end{lstlisting}

Raw traces go to \texttt{output/part3/}. The summaries, logs, and plots used
here are in \texttt{Assignment2/part3a-results/}.

\begin{itemize}
\item Experiment loop (Part 1 destinations, rotating order, retries):
  \href{\repo/part3/run_algorithms.py\#L224-L390}{\texttt{run\_part3}};
  per-run metrics:
  \href{\repo/part3/run_algorithms.py\#L126-L158}{\texttt{path\_metrics}};
  CLI options:
  \href{\repo/main.py\#L45-L96}{\texttt{main.py} L45--96}.
\item Algorithm selection and check:
  \href{\repo/part1/run_test.py\#L150-L170}{\texttt{run\_test.py} L150--170};
  goodput and loss proxy:
  \href{\repo/part1/tcp_stats.py\#L161-L257}{\texttt{compute\_sample}}.
\item Plots:
  \href{\repo/plot_results.py\#L601-L678}{\texttt{plot\_comparison}}
  (per-algorithm pages are also in \texttt{part3a-results/comparison.pdf}).
\end{itemize}

\section*{Use of AI tools}

We used Claude Code (Anthropic) to help set up the Linux VM, implement the
Part 3 changes, and draft this writeup from our measurement files. We
reviewed the code, tested it with short runs before the full experiment, and
checked the numbers against \texttt{part3a-results/}.

\end{document}
